\documentclass[11pt,a4paper]{article}
% preamble.tex — estilo común de todos los apuntes Froth.
% Cada apunte hace % preamble.tex — estilo común de todos los apuntes Froth.
% Cada apunte hace % preamble.tex — estilo común de todos los apuntes Froth.
% Cada apunte hace \input{preamble} justo después de \documentclass.
\usepackage[margin=2.4cm]{geometry}
\usepackage{xcolor}
\definecolor{litblue}{RGB}{31,79,121}
\definecolor{litgreen}{RGB}{27,94,32}
\definecolor{litorange}{RGB}{230,81,0}
\usepackage{parskip}
\usepackage{titlesec}
\titleformat{\section}{\Large\bfseries\color{litblue}}{\thesection}{0.6em}{}
\titleformat{\subsection}{\large\bfseries\color{litblue}}{\thesubsection}{0.6em}{}
\usepackage{tcolorbox}
\usepackage{fancyhdr}
\pagestyle{fancy}
\fancyhf{}
\fancyhead[L]{\small\color{litblue}Froth — Apuntes de ML}
\fancyhead[R]{\small\thepage}
\renewcommand{\headrulewidth}{0.4pt}
\usepackage[colorlinks=true,linkcolor=litblue,urlcolor=litblue]{hyperref}

% Cabecera de cada apunte: \cabecera{numero}{titulo}
\newcommand{\cabecera}[2]{%
  \begin{tcolorbox}[colback=litblue,coltext=white,colframe=litblue,arc=2mm]
    {\Large\bfseries Apunte #1 — #2}\\[3pt]
    {\small Proyecto Froth · aprender Machine Learning construyendo}
  \end{tcolorbox}\vspace{0.8em}}

% Cajas temáticas
\newtcolorbox{concepto}[1]{colback=blue!4,colframe=litblue,title={Concepto: #1},arc=1.5mm}
\newtcolorbox{practica}{colback=green!5,colframe=litgreen,title={Pruébalo tú (teclado en mano)},arc=1.5mm}
\newtcolorbox{ojo}{colback=orange!8,colframe=litorange,title={Ojo / error típico},arc=1.5mm}
\newtcolorbox{resumen}{colback=gray!8,colframe=gray!60,title={En una frase},arc=1.5mm}

\newcommand{\codigo}[1]{\texttt{#1}}
 justo después de \documentclass.
\usepackage[margin=2.4cm]{geometry}
\usepackage{xcolor}
\definecolor{litblue}{RGB}{31,79,121}
\definecolor{litgreen}{RGB}{27,94,32}
\definecolor{litorange}{RGB}{230,81,0}
\usepackage{parskip}
\usepackage{titlesec}
\titleformat{\section}{\Large\bfseries\color{litblue}}{\thesection}{0.6em}{}
\titleformat{\subsection}{\large\bfseries\color{litblue}}{\thesubsection}{0.6em}{}
\usepackage{tcolorbox}
\usepackage{fancyhdr}
\pagestyle{fancy}
\fancyhf{}
\fancyhead[L]{\small\color{litblue}Froth — Apuntes de ML}
\fancyhead[R]{\small\thepage}
\renewcommand{\headrulewidth}{0.4pt}
\usepackage[colorlinks=true,linkcolor=litblue,urlcolor=litblue]{hyperref}

% Cabecera de cada apunte: \cabecera{numero}{titulo}
\newcommand{\cabecera}[2]{%
  \begin{tcolorbox}[colback=litblue,coltext=white,colframe=litblue,arc=2mm]
    {\Large\bfseries Apunte #1 — #2}\\[3pt]
    {\small Proyecto Froth · aprender Machine Learning construyendo}
  \end{tcolorbox}\vspace{0.8em}}

% Cajas temáticas
\newtcolorbox{concepto}[1]{colback=blue!4,colframe=litblue,title={Concepto: #1},arc=1.5mm}
\newtcolorbox{practica}{colback=green!5,colframe=litgreen,title={Pruébalo tú (teclado en mano)},arc=1.5mm}
\newtcolorbox{ojo}{colback=orange!8,colframe=litorange,title={Ojo / error típico},arc=1.5mm}
\newtcolorbox{resumen}{colback=gray!8,colframe=gray!60,title={En una frase},arc=1.5mm}

\newcommand{\codigo}[1]{\texttt{#1}}
 justo después de \documentclass.
\usepackage[margin=2.4cm]{geometry}
\usepackage{xcolor}
\definecolor{litblue}{RGB}{31,79,121}
\definecolor{litgreen}{RGB}{27,94,32}
\definecolor{litorange}{RGB}{230,81,0}
\usepackage{parskip}
\usepackage{titlesec}
\titleformat{\section}{\Large\bfseries\color{litblue}}{\thesection}{0.6em}{}
\titleformat{\subsection}{\large\bfseries\color{litblue}}{\thesubsection}{0.6em}{}
\usepackage{tcolorbox}
\usepackage{fancyhdr}
\pagestyle{fancy}
\fancyhf{}
\fancyhead[L]{\small\color{litblue}Froth — Apuntes de ML}
\fancyhead[R]{\small\thepage}
\renewcommand{\headrulewidth}{0.4pt}
\usepackage[colorlinks=true,linkcolor=litblue,urlcolor=litblue]{hyperref}

% Cabecera de cada apunte: \cabecera{numero}{titulo}
\newcommand{\cabecera}[2]{%
  \begin{tcolorbox}[colback=litblue,coltext=white,colframe=litblue,arc=2mm]
    {\Large\bfseries Apunte #1 — #2}\\[3pt]
    {\small Proyecto Froth · aprender Machine Learning construyendo}
  \end{tcolorbox}\vspace{0.8em}}

% Cajas temáticas
\newtcolorbox{concepto}[1]{colback=blue!4,colframe=litblue,title={Concepto: #1},arc=1.5mm}
\newtcolorbox{practica}{colback=green!5,colframe=litgreen,title={Pruébalo tú (teclado en mano)},arc=1.5mm}
\newtcolorbox{ojo}{colback=orange!8,colframe=litorange,title={Ojo / error típico},arc=1.5mm}
\newtcolorbox{resumen}{colback=gray!8,colframe=gray!60,title={En una frase},arc=1.5mm}

\newcommand{\codigo}[1]{\texttt{#1}}

\begin{document}
\cabecera{18}{Git y GitHub: puntos de guardado y respaldo en la nube}

\section{Qué es Git}

\begin{concepto}{Git = máquina de puntos de guardado}
Como en un videojuego: cada vez que llegas a un estado que vale la pena conservar, haces un
\textbf{commit} (``guardar partida''): Git fotografía TODOS los archivos del proyecto en ese
instante y guarda la foto para siempre, con fecha, autor y un mensaje. Si mañana rompes algo,
puedes volver a cualquier foto anterior. El historial vive en la carpeta oculta \codigo{.git}.
Lo usa todo programador del planeta, a diario.
\end{concepto}

Vocabulario mínimo:
\begin{itemize}
  \item \textbf{repo(sitorio)}: el proyecto vigilado por Git.
  \item \codigo{git add .}: poner archivos ``en la bandeja'' para la próxima foto.
  \item \codigo{git commit -m "mensaje"}: disparar la foto.
  \item \codigo{git status}: ver qué hay en la bandeja / qué cambió.
  \item \textbf{rama (branch)}: una línea de historia. La principal se llama \codigo{main}.
  \item \textbf{remote / origin}: el apodo de la copia en la nube.
  \item \codigo{git push}: empujar tus fotos locales a la nube.
\end{itemize}

\begin{concepto}{.gitignore = lo que JAMÁS entra en las fotos}
Un archivo de texto con la lista de lo prohibido: secretos (\codigo{.openalex\_key}), la
burbuja (\codigo{.venv/}), datos regenerables (\codigo{2\_Datos/}), herramientas descargadas
(\codigo{tools/}). Así el repo lleva el \emph{código y el conocimiento}, no los pesos muertos
ni las llaves. Verificamos con \codigo{git ls-files} que ningún secreto entró: cero.
\end{concepto}

\section{Qué es GitHub (y qué NO es)}

\textbf{GitHub} es la nube donde respaldas los repos de Git — y además tu \textbf{portafolio
público}: cualquiera puede ver tu código, clonarlo (\codigo{git clone}) y correrlo en su PC.
GitHub \emph{guarda} el código; no lo \emph{ejecuta} por el visitante.

Tu repo: \codigo{github.com/kmortizva-data/froth} — de momento \textbf{privado} (GitHub
muestra 404 a los extraños para no revelar ni que existe). Cuando esté pulido:
Settings $\rightarrow$ Change visibility $\rightarrow$ Public, y nace el portafolio.

\section{Lo que pasó en tu primera vez (incluye los tropiezos, que enseñan)}

\begin{enumerate}
  \item \codigo{git init} + \codigo{git add .} + \codigo{git commit} $\rightarrow$ primera
        foto: 65 archivos, cero secretos.
  \item \textbf{Tropiezo 1}: el push falló con ``repository does not seem to exist'' — el
        remoto apuntaba a una URL de ejemplo cuya página nunca se creó en GitHub. Lección:
        \emph{el repo en la nube hay que crearlo antes de empujar} (o dejar que GitHub
        Desktop lo cree con su botón \textbf{Publish repository}, que hace ambas cosas).
  \item \textbf{Tropiezo 2}: \codigo{git remote remove origin} respondió ``No such remote''
        — significaba ``ya estaba borrado'', no un error real. Lección: lee el mensaje;
        a veces el ``error'' es solo un aviso de que no había nada que hacer.
  \item \codigo{git branch -M main}: renombrar \codigo{master} $\rightarrow$ \codigo{main}
        (el estándar moderno). El silencio de la terminal $=$ éxito.
  \item \textbf{Publish repository} en GitHub Desktop $\rightarrow$ creó el repo en la nube y
        empujó. Verificado: \codigo{main...origin/main} sin diferencias $=$ sincronizado.
\end{enumerate}

\section{Tu rutina desde ahora (la única que necesitas memorizar)}

Al final de cada sesión de trabajo con cambios que valgan:
\begin{center}
\codigo{git add .} \quad$\rightarrow$\quad \codigo{git status} (revisar) \quad$\rightarrow$\quad
\codigo{git commit -m "qué hice"} \quad$\rightarrow$\quad \codigo{git push}
\end{center}
(O en GitHub Desktop: revisar cambios $\rightarrow$ escribir resumen $\rightarrow$ Commit
$\rightarrow$ Push. Es lo mismo con botones.)

\begin{resumen}
Git guarda fotos completas del proyecto (commits) y GitHub las respalda en la nube y las
exhibe como portafolio; el .gitignore protege secretos y pesos muertos. Tu rutina: add
$\rightarrow$ status $\rightarrow$ commit $\rightarrow$ push. Primera foto: 65 archivos,
cero secretos, ya en la nube.
\end{resumen}

\end{document}
